\documentclass{article}
\usepackage{amsmath}
\usepackage{amssymb}
\usepackage{array}
\usepackage{booktabs}
\usepackage{textcomp}
\usepackage{graphicx}

\newcommand{\coNP}{\textrm{co-}\mathbb{NP}}

\title{hw1}
\date{\today}
\begin{document}
\maketitle

\section{ChatGPT questions} 

Put this question ChatGPT and get the answer. You are allowed to discuss further. Please put a link to the response and your conversation. 
Given the response, you must first say whether: (i) the response is correct, (ii) the response is false, or (iii) you are really not sure.
In the first case, give a succinct explanation of the proof. In the second case, explain why it is false.
In the third case, explain your confusion. You do not need to write too much.

Conversations can be found here: 

\begin{enumerate}
    \item Give a bijection from $\{x | x > 1, x \in \mathbb R\}$ and the interval $(0,1)$.
    

    \textbf{claude 3.5 response}

    I'll solve this by constructing a bijective function between the set of real numbers greater than 1 and the open interval (0,1).

    \begin{enumerate}
        \item Bijection Construction
    
        A classic and elegant solution to this problem is to use the following function:
    
        $f(x) = \frac{1}{x-1}$
    
        \item Proof of Bijectivity
    
        \begin{enumerate}
            \item Injective (One-to-One)
            Let $x_1, x_2 \in \{x | x > 1, x \in \mathbb R\}$
            If $f(x_1) = f(x_2)$, then:
            $\frac{1}{x_1-1} = \frac{1}{x_2-1}$
            This implies $x_1 = x_2$, proving injectivity.
        
            \item Surjective (Onto)
            For any $y \in (0,1)$, we can find an $x$ such that $f(x) = y$:
            $y = \frac{1}{x-1}$
            $x = \frac{1}{y} + 1$
        \end{enumerate}
    
        Note that:
        - When $y \in (0,1)$, $x = \frac{1}{y} + 1 > 1$
        - $f(x) = \frac{1}{(\frac{1}{y} + 1) - 1} = y$
    
        Thus, $f$ maps $\{x | x > 1, x \in \mathbb R\}$ bijectively onto $(0,1)$.




    \end{enumerate}
    \textbf{End response}

    the response is correct. it shows that for \(x>1, x\in \mathbb R\), the map \(f(x)=\frac{1}{x-1}\) is injective and surjective to the interval \((0,1)\). 




\end{enumerate}
\begin{enumerate}
    \item NAE-3SAT is the language of 3-CNFs with an assignment that does not make all literals in a clause have the same value.
    Prove that NAE-3SAT is $\mathbb{NP}$-complete. 

    \textbf{response}

    the proof is incorrect. the mapping from 3sat to nae3sat is invalid since the clauses are not 3 literals, and it doesnt give a way to convert the 4 literals into 3 literals.



    \item Prove that if FACTORING is $\coNP$-complete, then $TAUT \leq_{p(n)} SAT$ (recall
    that $TAUT$ is the language of tautologically true Boolean formulas.)

    \textbf{response}

    i am not sure if this proof is correct. it is almost certainly wrong but the words dont make any sense. claude is totally hallucinating. ill try it with chatgpt


    

\end{enumerate}

\end{document}